Controllers come in a multitude of flavors. Some are designed to withstand gradual faults, while others are designed to stabilize a \gls{S/C} after a total thruster failure. Worth mentioning before this section is that no controller, whether it be \gls{SMC}, \gls{MPC}, or a PID, is inherently active/passive, nor are they robust, adaptive, or intelligent. The terms come from the system structure around the controller.

\subsection{Robust Control}

\gls{RC} refers to control methods designed to maintain stability and acceptable performance in the presence of bounded uncertainties and disturbances \cite{Zhang2008BibliographicalSystems,Jiang2012Fault-tolerantApproaches}. Unlike adaptive schemes, robust controllers are fixed once implemented; they rely on conservative design margins rather than online adjustment \cite{Jiang2012Fault-tolerantApproaches}.

Classical examples include $H_\infty$ control, which minimizes the worst-case disturbance amplification, as well as $\mu$-synthesis and reliable control formulations \cite{Yin2016ASystem}. These approaches are especially attractive when actuator faults or parameter variations can be bounded in advance, allowing the controller to absorb the effects without reconfiguration. For spacecraft, robust designs based on $H_\infty$ have been proposed for attitude stabilization under disturbances and partial actuator degradations, providing stability guarantees without requiring explicit fault diagnosis \cite{Yin2016ASystem}.

The main limitation of robust control is that performance is guaranteed only within the predefined uncertainty set. Larger or unexpected faults may exceed the controller’s robustness margins, leading to degraded performance or instability. For this reason, robust controllers are often used as passive FTC baselines, or combined with reconfiguration mechanisms in an active FTC framework to expand their fault coverage \cite{Zhang2008BibliographicalSystems,Jiang2012Fault-tolerantApproaches}.

\subsection{Sliding Mode Control (SMC)}
  \gls{SMC} is a nonlinear control approach that forces the system trajectory onto a predefined “sliding surface,” where the reduced-order dynamics are stable and robust to matched uncertainties \cite{Yin2016ASystem,Mirshams2014PassiveThrusters}. Its strength lies in handling abrupt actuator faults and external disturbances without requiring an exact model of the spacecraft dynamics. Variants such as terminal SMC and fixed-time SMC guarantee convergence to the desired state within bounded time \cite{Cao2013FaultMode}, while velocity-free SMC formulations eliminate the need for direct angular-rate measurements by employing state observers \cite{Guo2019Velocity-freeSaturation}
.

In spacecraft applications, SMC has been demonstrated as particularly effective for thruster-actuated platforms. For example, Mirshams et al. \cite{Mirshams2014PassiveThrusters} showed that an SMC-based design could maintain attitude stability after discrete thruster failures, leveraging the inherent robustness of the sliding surface dynamics. These properties make SMC attractive for \gls{FTC} scenarios where actuators fail in an on/off manner, a common fault mode in spacecraft propulsion systems.

However, SMC also carries limitations. The well-known chattering effect—caused by high-frequency switching—can excite unmodeled dynamics and degrade performance \cite{Yang2023AdaptiveFaults}. Moreover, SMC does not naturally account for input constraints such as actuator saturation, which are ubiquitous in spacecraft. To address these shortcomings, recent research explores hybrid approaches that combine SMC with optimization-based allocation or predictive control, enabling constraint handling while preserving robustness \cite{Khodaverdian2023Fault-tolerantSpacecraft}
.


\subsection{Adaptive Control}
\gls{AC} methods adjust controller parameters online to compensate for changing dynamics or actuator effectiveness. In the context of \gls{FTC}, they are particularly suited for gradual or parametric faults, such as partial loss of actuator authority or slow variations in inertia \cite{Jiang2012Fault-tolerantApproaches,Yin2016ASystem}. Rather than relying on fixed gains, an adaptive law continuously estimates system changes and updates the controller in real time.

Common spacecraft implementations include \\
\gls{MRAC}, which enforces tracking of a predefined reference model despite actuator failures \cite{Ma2014ActuatorFeedback}, and adaptive backstepping, which provides stability guarantees while compensating for parametric uncertainties \cite{Bustan2014AdaptiveControl}. Adaptive extensions of other families, such as adaptive SMC \cite{Yang2023AdaptiveFaults} and adaptive control allocation \cite{Tohidi2016FaultSpace}, have also been proposed to broaden fault tolerance while preserving robustness.

The strength of adaptive control lies in its ability to recover performance smoothly without explicit controller switching or redesign \cite{Jiang2012Fault-tolerantApproaches}. However, this benefit comes at the cost of slower response to abrupt, discrete faults, since adaptation requires time to estimate and converge to new parameters \cite{Zhang2008BibliographicalSystems}. Furthermore, the stability and effectiveness of adaptive FTC depend strongly on the accuracy of the fault estimation and the tuning of the adaptation law. Poor estimation quality can lead to oscillations or even instability, which has motivated hybrid designs that combine adaptive controllers with robust supervisory layers \cite{Zhang2008BibliographicalSystems,Jiang2012Fault-tolerantApproaches}.

\subsection{Intelligent, Fuzzy, and Learning-Based Control} 
\gls{FTC} methods increasingly explore intelligent and data-driven approaches that complement or replace model-based designs with heuristic or learning-based adaptation. \gls{FLC}s use linguistic rules (“if actuator efficiency decreases, then increase control effort”) to compensate for faults without requiring precise mathematical models. This makes them effective for gradual or ambiguous degradations, where exact fault dynamics are uncertain \cite{Zhang2008BibliographicalSystems,Yin2016ASystem}. Adaptive fuzzy controllers have been applied to spacecraft attitude control, demonstrating resilience to actuator faults by updating rule weights online\cite{Zou2011AdaptiveSpacecraft}
.

Beyond fuzzy logic, \gls{NN}s and \gls{RL} extend the paradigm by learning nonlinear mappings or control policies from data. \gls{NN}s can approximate unmodeled actuator or sensor fault effects, while \gls{RL} agents can directly learn fault-tolerant policies through trial-and-error optimization. For spacecraft, these approaches are promising for complex and poorly modeled faults such as partial thruster degradations or sensor bias. Recent work has demonstrated RL-based fault-tolerant attitude control, where an agent learns to stabilize a satellite despite actuator faults \cite{Rafiee2024ActiveFaults}.

While intelligent controllers offer flexibility, they remain less flight-proven than robust or adaptive methods. Stability guarantees are weaker, and their certification for space missions is challenging. As a result, fuzzy, neural, and RL-based approaches are most commonly found in simulations or laboratory testbeds rather than operational spacecraft \cite{Zhang2008BibliographicalSystems,Yin2016ASystem}. Ongoing research therefore focuses on hybrid strategies that combine learning-based adaptation with robust supervisory control, aiming to merge the flexibility of intelligent methods with the assurance required for safety-critical spacecraft operations.